\documentclass[11pt]{article}
\usepackage[T1]{fontenc}
\usepackage{amsmath,amssymb,amsthm}
\usepackage[margin=1in]{geometry}
\usepackage{hyperref}
\hypersetup{
  pdftitle={Exact Condition-Number Asymptotics for Four Matrices in Sylvester Example 6},
  pdfauthor={Xiaojian Zeng}
}

\newtheorem{theorem}{Theorem}
\newtheorem{lemma}[theorem]{Lemma}
\newcommand{\R}{\mathbb{R}}
\newcommand{\norm}[1]{\lVert#1\rVert_2}

\title{Exact Condition-Number Asymptotics for Four Matrices in Sylvester Example~6}
\author{Xiaojian Zeng\\
PBC School of Finance, Tsinghua University, 100190\\
\texttt{xiaojianzeng@foxmail.com}}
\date{}

\begin{document}
\maketitle

\begin{abstract}
For four right-hand sides in Example~6 of Fasi and Hashemi's study of
Sylvester-equation solutions, numerical experiments suggested condition-number
orders $\delta^{-2}$, $\delta^{-1}$, $1$, and $1$.  We determine the exact
leading constants: respectively
\[
27,\qquad
\sqrt{\frac{778752}{57905-5\sqrt{132562057}}},\qquad
\sqrt{\rho_3/\rho_1},\qquad \frac{13}{3},
\]
where $\rho_1<\rho_2<\rho_3$ are explicitly isolated roots of a stated
integer cubic.  The two bounded limits are thus obtained exactly.  We also
give explicit linear relative-error bounds, including full-interval bounds in
the first and fourth cases.  All four matrices are nonsingular throughout the
parameter interval, so no rank-deficient convention is needed.
\end{abstract}

\section{Setting and result}

Let $0<\delta\leq 1$ and
\[
A=\operatorname{diag}(1,4,7),\qquad
B_\delta=\operatorname{diag}(10,1-\delta,20).
\]
For $C_1=I+E_{12}$, $C_2=\mathbf 1$, $C_3=\mathbf 1-E_{12}$, and
$C_4=I$, let $X_i(\delta)$ solve
\[
 AX_i-X_iB_\delta=C_i.
\]
The spectra of $A$ and $B_\delta$ are disjoint, so the solution is unique.
Entrywise solution of the equation gives
\[
X_1=\begin{pmatrix}
-1/9&1/\delta&0\\0&1/(3+\delta)&0\\0&0&-1/13
\end{pmatrix},
\quad
X_2=\begin{pmatrix}
-1/9&1/\delta&-1/19\\
-1/6&1/(3+\delta)&-1/16\\
-1/3&1/(6+\delta)&-1/13
\end{pmatrix},
\]
\[
X_3=\begin{pmatrix}
-1/9&0&-1/19\\
-1/6&1/(3+\delta)&-1/16\\
-1/3&1/(6+\delta)&-1/13
\end{pmatrix},
\quad
X_4=\operatorname{diag}\!\left(-\frac19,\frac1{3+\delta},-\frac1{13}\right).
\tag{1}
\]
For a full-rank matrix $Y$, write
$\kappa_2(Y)=\sigma_1(Y)/\sigma_3(Y)$, with singular values in decreasing
order.

Fasi and Hashemi~\cite{FasiHashemi2026} plotted these four condition numbers
at 30 parameter values and reported the four orders above, without exact
constants or remainder estimates.  Exact algebraic constants supported by
finite-dimensional matrix certificates also occur in the distinct
spectral-set problem studied by Zeng~\cite{Zeng2026ConvexLens}; that work does
not address Sylvester equations.  Our contribution here is confined to the
four matrices in (1).  We use the standard singular-value perturbation bound
$|\sigma_j(Y+E)-\sigma_j(Y)|\leq\norm{E}$; see, for example,
Stewart and Sun~\cite[Theorem~IV.4.11]{StewartSun1990}.

Define
\[
\mu_-:=\frac{57905-5\sqrt{132562057}}{778752},\qquad s:=\sqrt{\mu_-},
\tag{2}
\]
and let $0<\rho_1<\rho_2<\rho_3$ be the roots of
\[
P(z)=45542974464z^3-13786630308z^2+429202697z-330625,
\tag{3}
\]
uniquely specified by
\[
\rho_1\in(1/1300,1/1200),\quad
\rho_2\in(1/30,1/29),\quad
\rho_3\in(1/4,1/3).
\tag{4}
\]
For the explicit error constants in the second case put
\[
U=\frac{\sqrt{442}}{171},\quad V=\frac{\sqrt5}{6},\quad
a=UV=\frac{\sqrt{2210}}{1026},
\]
\[
c=U+V+a,\quad F=2c+c^2,\quad b=a/s,
\quad M_2=F+2b+2Fb,\quad
\delta_2=\min\{1,s/(2a)\}.
\tag{5}
\]
For the third case put
\[
e=\frac{\sqrt{17}}{36},\quad
\alpha=\frac{e}{\sqrt{\rho_3}},\quad
\beta=\frac{e}{\sqrt{\rho_1}},\quad
M_3=\alpha+2\beta+2\alpha\beta.
\tag{6}
\]

\begin{theorem}
Every $X_i(\delta)$ in (1) is nonsingular for $0<\delta\leq1$.  As
$\delta\to0^+$,
\[
\kappa_2(X_1)\sim27\delta^{-2},\qquad
\kappa_2(X_2)\sim
\sqrt{\frac{778752}{57905-5\sqrt{132562057}}}\,\delta^{-1},
\tag{7}
\]
\[
\kappa_2(X_3)\longrightarrow\sqrt{\rho_3/\rho_1},
\qquad
\kappa_2(X_4)\longrightarrow\frac{13}{3}.
\tag{8}
\]
More precisely,
\[
\left|\frac{\kappa_2(X_1)}{27\delta^{-2}}-1\right|
 \leq\frac{121}{243}\delta
\quad(0<\delta\leq1),
\tag{9}
\]
\[
\left|\frac{\kappa_2(X_2)}{s^{-1}\delta^{-1}}-1\right|
 \leq M_2\delta
\quad(0<\delta\leq\delta_2),
\tag{10}
\]
\[
\left|\frac{\kappa_2(X_3)}{\sqrt{\rho_3/\rho_1}}-1\right|
 \leq M_3\delta
\quad(0<\delta\leq1/10),
\tag{11}
\]
and, throughout $0<\delta\leq1$,
\[
\kappa_2(X_4)=\frac{13}{3+\delta},\qquad
\left|\frac{\kappa_2(X_4)}{13/3}-1\right|\leq\frac\delta3.
\tag{12}
\]
\end{theorem}

\section{Proof}

Direct determinant expansion gives
\[
\det X_1=\det X_4=\frac1{117(3+\delta)},
\]
\[
\det X_2=\frac{5(\delta+9)(\delta+19)}
 {5928\delta(\delta+3)(\delta+6)},\qquad
\det X_3=-\frac{5(17\delta+115)}
 {11856(\delta+3)(\delta+6)}.
\tag{13}
\]
None vanishes on $(0,1]$, proving the full-rank assertion.

\subsection*{The first matrix}

The nontrivial $2\times2$ block of $X_1$ is
\[
T=\begin{pmatrix}-1/9&1/\delta\\0&1/(3+\delta)\end{pmatrix}.
\]
Let $\lambda_+\geq\lambda_->0$ be the eigenvalues of $T^TT$.  Rayleigh
quotients and the trace give
\[
\delta^{-2}\leq\lambda_+\leq
\delta^{-2}+\frac1{81}+\frac1{(3+\delta)^2}
\leq\delta^{-2}+\frac{10}{81}.
\tag{14}
\]
Since $\lambda_+\lambda_-=1/[81(3+\delta)^2]$,
\[
\lambda_-\leq\frac{\delta^2}{81(3+\delta)^2}
\leq\frac1{1296}<\frac1{169}.
\tag{15}
\]
Thus the largest and smallest squared singular values of $X_1$ are
$\lambda_+$ and $\lambda_-$, rather than the scalar value $1/169$.
Consequently
\[
\kappa_2(X_1)=\sqrt{\frac{\lambda_+}{\lambda_-}}
=9(3+\delta)\lambda_+.
\]
Equations (14)--(15) yield
\[
1+\frac\delta3\leq
\frac{\kappa_2(X_1)}{27\delta^{-2}}
\leq\left(1+\frac\delta3\right)
       \left(1+\frac{10\delta^2}{81}\right).
\tag{16}
\]
After subtracting one, the right side of (16) is at most
$121\delta/243$ for $0<\delta\leq1$.  This proves (7) and (9) for $X_1$.

\subsection*{The second matrix}

Permuting the columns of $X_2$ in the order $(2,1,3)$ preserves all singular
values and gives
\[
M=\begin{pmatrix}t&u^T\\v&D\end{pmatrix},\quad t=\delta^{-1},\quad
u=\begin{pmatrix}-1/9\\-1/19\end{pmatrix},\quad
v=\begin{pmatrix}1/(3+\delta)\\1/(6+\delta)\end{pmatrix},
\]
\[
D=\begin{pmatrix}-1/6&-1/16\\-1/3&-1/13\end{pmatrix}.
\tag{17}
\]
Block multiplication verifies the exact factorization
\[
M=
\begin{pmatrix}1&0\\\delta v&I_2\end{pmatrix}
\begin{pmatrix}t&0\\0&D-\delta vu^T\end{pmatrix}
\begin{pmatrix}1&\delta u^T\\0&I_2\end{pmatrix}.
\tag{18}
\]
Here $\norm{u}=U$ and $\norm{v}\leq V$.  Moreover,
\[
\det(zI_2-D^TD)=z^2-\frac{57905}{389376}z+
\frac{25}{389376},
\tag{19}
\]
so (2) is positive and $s=\sigma_{\min}(D)$.  The perturbation bound gives
\[
s-a\delta\leq\sigma_{\min}(D-\delta vu^T)
\leq s+a\delta.
\tag{20}
\]

Each outer shear in (18), as well as its inverse, has norm at most
$1+\delta V$ or $1+\delta U$.  Set
$f=(1+\delta U)(1+\delta V)$.  Multiplication by the two shears therefore
changes each singular value of the middle block by a factor in
$[f^{-1},f]$.  Also
\[
\norm{D-\delta vu^T}\leq\norm{D}+a<1\leq t;
\]
the strict inequality follows, for example, from
$\norm{D}\leq1/6+1/16+1/3+1/13<2/3$ and $a<1/20$.
For $\delta\leq\delta_2$, (20) is positive, so the largest and smallest
singular values of the middle block are $t$ and
$\sigma_{\min}(D-\delta vu^T)$.  It follows that
\[
\frac1{f^2}\frac{s}{s+a\delta}
\leq\frac{\kappa_2(X_2)}{s^{-1}\delta^{-1}}
\leq f^2\frac{s}{s-a\delta}.
\tag{21}
\]
For $0<\delta\leq1$,
$f\leq1+c\delta$ and hence $f^2\leq1+F\delta$.  Since
$b\delta\leq1/2$ on $(0,\delta_2]$, expanding the two rational bounds in
(21) gives the stated coefficient below; direct comparison of the squares in
(2) and (5) gives $a>s$, hence $b>1/2$:
\[
\left|\frac{\kappa_2(X_2)}{s^{-1}\delta^{-1}}-1\right|
\leq(F+2b+2Fb)\delta.
\]
Together with (2), this is (7) and (10).

\subsection*{The third matrix}

Let $X_{3,0}=X_3(0)$, where the notation means continuous substitution in
the bounded entries of (1).  Exact expansion gives
\[
\det(zI_3-X_{3,0}^TX_{3,0})=\frac{P(z)}{45542974464}.
\tag{22}
\]
For completeness, the signs of the exact Sturm chain of $P$ at the endpoints
in (4), together with the variation counts, are
\[
\begin{array}{c|c|c|c}
\text{interval}&\text{left signs}&\text{right signs}&\text{variations}\\ \hline
(1/1300,1/1200)&(-,+,-,+)&(+,+,-,+)&3\to2\\
(1/30,1/29)&(+,-,+,+)&(-,-,+,+)&2\to1\\
(1/4,1/3)&(-,+,+,+)&(+,+,+,+)&1\to0.
\end{array}
\tag{23}
\]
Sturm's theorem gives exactly one root in each interval.  Since $P$ is cubic,
these are all its roots, and (4) proves their positivity and ordering.  Thus
the limiting squared singular values are $\rho_1<\rho_2<\rho_3$, proving the
limit in (8).

Only entries $(2,2)$ and $(3,2)$ vary with $\delta$, and direct subtraction
gives
\[
\norm{X_3(\delta)-X_{3,0}}
\leq\lVert X_3(\delta)-X_{3,0}\rVert_F
\leq e\delta.
\tag{24}
\]
The singular-value perturbation bound now gives
\[
\sqrt{\rho_3}-e\delta\leq\sigma_1(X_3(\delta))
\leq\sqrt{\rho_3}+e\delta,
\]
\[
\sqrt{\rho_1}-e\delta\leq\sigma_3(X_3(\delta))
\leq\sqrt{\rho_1}+e\delta.
\tag{25}
\]
Because $\rho_1>1/1300$, one has $\beta\delta<1/2$ for
$0<\delta\leq1/10$.  Dividing (25) by the limiting extreme singular values
therefore bounds the normalized condition number between
$(1-\alpha\delta)/(1+\beta\delta)$ and
$(1+\alpha\delta)/(1-\beta\delta)$.  Elementary expansion gives (11): the
stated $M_3$ dominates the two resulting coefficients (in particular,
$\beta>1/2$).  This completes the third case.

\subsection*{The fourth matrix}

For $0<\delta\leq1$, the largest magnitude among the three diagonal entries
of $X_4$ is $1/(3+\delta)$ and the smallest is $1/13$.  Hence (12) follows
immediately, completing the proof.

\section*{Computational and AI assistance}

An AI language model assisted with algebraic derivation and manuscript
drafting.  During preparation, a deterministic SymPy script using exact
rational and symbolic arithmetic independently expanded the four Gram-matrix
characteristic polynomials, checked (13), (19), and (22), and verified the
three Sturm root counts in (23).  Decimal singular-value evaluations were used
only as diagnostics and play no role in the proof.

\bibliographystyle{plain}
\bibliography{references}

\end{document}
